\documentclass[10pt,a4paper]{amsart}
\usepackage[margin=19mm]{geometry}
\usepackage{amsmath,amssymb,mathtools}
\usepackage[scr=rsfs,cal=euler]{mathalfa}
\usepackage{xfrac}
\usepackage[unicode,hidelinks,bookmarks=false]{hyperref}
\newcommand{\ie}{i.e.\ }
\newcommand{\cf}{cf.\ }
\newcommand{\resp}{respectively\ }
\newcommand{\Subsection}{\subsection}
\providecommand{\nobreakdash}{\nobreak\hbox{-}}
\setlength{\emergencystretch}{2em}
\multlinegap0pt
\usepackage{amssymb}
\usepackage{enumerate}
\usepackage{float}
\usepackage{cancel}
\usepackage{mathtools}
\newtheorem{lemma}{Lemma}
\newtheorem{case}{Case}
\newcommand{\M}{\mathcal{M}}
\newcommand{\R}{{\mathbb R}}
\DeclareMathOperator{\arccot}{\mathrm{arccot}}
\DeclareMathOperator{\area}{\mathrm{area}}
\newcommand{\hyp}{{\mathbb H}}
\newcommand{\sph}{{\mathbb S}}
\title{Optimal stability of P\'al's isominwidth inequality for ball convex bodies in planes of constant curvature}
\author{Ferenc Fodor, \'Ad\'am Sagmeister}
\date{Source: arXiv:2606.02882v1 (CC BY 4.0)}
\begin{document}
\maketitle

\noindent\emph{Source and licence.} Optimal stability of P\'al's isominwidth inequality for ball convex bodies in planes of constant curvature. Version arXiv:2606.02882v1 (CC BY 4.0), distributed by arXiv under the Creative Commons Attribution 4.0 International licence, http://creativecommons.org/licenses/by/4.0/, as recorded in the arXiv licence metadata for this exact version and shown on the abstract page https://arxiv.org/abs/2606.02882v1. Full licence notice: \texttt{licenses/arxiv-2606.02882-CC-BY-4.0.txt}.\par\smallskip

\noindent\textit{Editorial scope.} Counted unit: the article's lemma lem:distance-of-K-and-C (source0.tex lines 467--473). The article's complete original proof of it is delivered with it (source0.tex lines 475--541, one \texttt{proof} environment of 67 lines). No mathematics is written, repaired or re-derived: the statement and the proof are the article's own lines, in the article's own notation, and no retained line is substituted or re-flowed.  Retained uncounted as the source-local context the proof needs: lines 280--284.  Nothing was omitted from any retained span, so the package carries no omission record.  No retained line contains a citation, so the wrapper carries no bibliography and no bibliography entry.  No retained line contains a figure environment, an image-inclusion command or drawing code, so no image is delivered and the package declares no asset dependency.  Editorial formatting only, in addition: an AMS \texttt{amsart} preamble that restates the article's own declarations and macros used by the retained lines, namely the environments \texttt{lemma}, \texttt{case} on separate counters, together with the standard packages those lines need.  The retained environments print in this document as Lemma 1, Case 1 in order of appearance, whereas the article prints the same items with the numbering of its own full document.  The complete native source of the article as distributed in the arXiv e-print for this exact version is bundled unchanged as \texttt{sources/201935/source0.tex}.
\begin{lemma}\label{lemma:triangle-area}
Let $T\subset\M^2$ be a triangle of side lengths $a,b,c$ and opposite angles $\alpha,\beta,\gamma$. Then for the area of $T$ we have the following equations depending on $\M^2$:
$$\cot\left(\frac{\area\left(T\right)}{2}\right)=\frac{\coth\frac{a}{2}\coth\frac{b}{2}-\cos\gamma}{\sin\gamma}\text{ if }\M^2=\hyp^2,$$
$$\cot\left(\frac{\area\left(T\right)}{2}\right)=\frac{\cot\frac{a}{2}\cot\frac{b}{2}+\cos\gamma}{\sin\gamma}\text{ if }\M^2=\sph^2.$$
\end{lemma}

\begin{lemma}\label{lem:distance-of-K-and-C}
Let $K$ be an $r$-ball convex body of minimal width $0<w\leq r$ and $\area(K)\leq\area(T_{w,r})+\varepsilon$ with $\varepsilon\in(0,\overline{\varepsilon})$. Let $C=B\cup C_1\cup C_2\cup C_3$ be a cap-domain with non-overlapping congruent caps $C_i$ peaks at $q_i$, where $d(q_i,p)=w-\varrho_0-\eta$. Then
$$
\delta(K,C)<c_4\varepsilon
$$
for some positive constant $c_4$.
\end{lemma}

\begin{proof}
Let $K_1,K_2,K_3$ be the regions containing $C_1,C_2,C_3$, respectively, bounded by circular arcs of $\partial B$ and parts of $\partial K$. Let $z$ be a farthest point of $K$ from $p$. We may assume that $z\in K_1$. Moreover, by the ball convexity of $K$, we may assume $q_1\in[p,z]$. Then, it is enough to show that $d(z,q_1)\leq c_4\varepsilon$ for some positive constant $c_4$. If $q_1=z$, then there is nothing to prove, so we assume otherwise. Since $w-\varrho_0-\eta>\varrho_0+\eta$, $d(q_1,s_{1,2})>d(p,s_{1,2})>\varrho_0$ and $\angle(s_{1,3},q_1,s_{1,2})<\angle(s_{1,3},p,s_{1,2})<\frac{2\pi}{3}$. By our assumption $q_1\in[p,z]$, $\angle(s_{1,3},q_1,s_{1,2})+\angle(s_{1,3},q_1,z)=\angle(s_{1,3},q_1,s_{1,2})+\angle(s_{1,2},q_1,z)>\pi$ and $\angle(z,q_1,s_{1,2})=\angle(z,q_1,s_{1,3})>\frac{2\pi}{3}$. On the other hand, $\angle(q_1,z,s_{1,2})<\angle(s_{1,3},z,s_{1,2})\leq\angle(s_{1,3},q_1,s_{1,2})<\frac{2\pi}{3}$. This implies
\begin{equation}\label{eq:sides-of-triangle-zs12q1}
d(s_{1,2},q_1)<d(z,s_{1,2}).
\end{equation}
Since $K$ is $r$-ball convex, the $r$-segment determined by $z$ and $s_{1,2}$ (the $r$-ball convex hull of $z$ and $s_{1,2}$) is contained in $K$. In other words, the region $T$ bounded by $[q_1,z]$, a convex $r$-circular arc connecting $z$ and $s_{1,2}$, and a concave $r$-circular arc connecting $s_{1,2}$ and $q_1$ is contained in $K_1$. From \eqref{eq:sides-of-triangle-zs12q1} and our assumption,
\begin{equation}
\area([z,s_{1,2},q_1])<\area(T)<\varepsilon.
\end{equation}
The rest of the proof is divided into cases according to $\M^2$.

\setcounter{case}{0}
\begin{case}
    $\M^2=\R^2$.
\end{case}

In this case,
\[
\varepsilon>\area([z,s_{1,2},q_1])=\frac{d(q_1,s_{1,2})\cdot d(q_1,z)\sin(\angle(s_{1,2},q_1,z))}{2}>\frac{\sqrt{3}}{4}\varrho_0 d(q_1,z),
\]
and hence
\[
d(q_1,z)<\frac{4\sqrt{3}}{3\varrho_0}\varepsilon.
\]

\begin{case}
    $\M^2=\hyp^2$.
\end{case}

From the area formula of the hyperbolic triangle in Lemma \ref{lemma:triangle-area},
\[
\varepsilon>2\arccot\left(\frac{\coth\frac{d(s_{1,2},q_1)}{2}\coth\frac{d(z,q_1)}{2}-\cos\angle(s_{1,2},q_1,z)}{\sin\angle(s_{1,2},q_1,z)}\right)>2\arccot\left(\frac{\coth\frac{\varrho_0}{2}\coth\frac{d(z,q_1)}{2}+1}{\frac{\sqrt{3}}{2}}\right),
\]
which is equivalent to
\[
\tanh\frac{\varrho_0}{2}\left(\frac{\sqrt{3}}{2}\cot\frac{\varepsilon}{2}-1\right)<\coth\frac{d(z,q_1)}{2}.
\]
We can bound the left-hand side from below, where we also use $\varepsilon\in\left(0,\frac{3}{4}\right)$.
\[
\tanh\frac{\varrho_0}{2}\left(\frac{\sqrt{3}}{2}\cot\frac{\varepsilon}{2}-1\right)>\tanh\frac{\varrho_0}{2}\coth(2\varepsilon)>\coth\left(2\coth\frac{\varrho_0}{2}\varepsilon\right),
\]
and thus
\[
d(q_1,z)<4\coth\frac{\varrho_0}{2}\varepsilon.
\]

\begin{case}
    $\M^2=\sph^2$.
\end{case}

From the area formula in Lemma \ref{lemma:triangle-area}, we have
\[
\varepsilon>2\arccot\left(\frac{\cot\frac{d(s_{1,2},q_1)}{2}\cot\frac{d(z,q_1)}{2}+\cos\angle(s_{1,2},q_1,z)}{\sin\angle(s_{1,2},q_1,z)}\right)>2\arccot\left(\frac{\cot\frac{\varrho_0}{2}\cot\frac{d(z,q_1)}{2}-\frac{1}{2}}{\frac{\sqrt{3}}{2}}\right),
\]
or equivalently
\[
\tan\frac{\varrho_0}{2}\left(\frac{\sqrt{3}}{2}\cot\frac{\varepsilon}{2}+\frac{1}{2}\right)<\cot\frac{d(z,q_1)}{2}.
\]
We can bound the left-hand side from below, using $\varepsilon\in(0,1)$.
\[
\tan\frac{\varrho_0}{2}\left(\frac{\sqrt{3}}{2}\cot\frac{\varepsilon}{2}+\frac{1}{2}\right)>\tan\frac{\varrho_0}{2}\cot\varepsilon>\cot\left(\cot\frac{\varrho_0}{2}\varepsilon\right),
\]
and thus
\[
d(q_1,z)<2\cot\frac{\varrho_0}{2}\varepsilon.
\]
\end{proof}

\end{document}
